\section{Source Code Snippets}

This appendix contains key source code implementations demonstrating core NexusFlow functionality across Android and firmware layers.

\subsection{Android: Command Execution Funnel}

\begin{lstlisting}[language=Kotlin, caption=Unified Relay Command Funnel, label=lst:android_command_funnel]
/**
 * Unified command funnel - all relay control paths converge here
 * Ensures consistent state management and auditability
 */
suspend fun applyCommand(
    deviceId: String,
    relayIndex: Int,
    action: RelayAction,
    source: CommandSource,
    metadata: CommandMetadata? = null
): Result<CommandResponse> = withContext(Dispatchers.Default) {
    try {
        // Pre-execution validation
        if (relayIndex !in 0..5) {
            return@withContext Result.failure(
                InvalidRelayException("Relay index $relayIndex out of range")
            )
        }
        
        val device = deviceRepository.getDevice(deviceId)
            ?: return@withContext Result.failure(
                DeviceNotFoundException("Device $deviceId not found")
            )
        
        // Check for conflicts
        if (source != CommandSource.AUTOMATION) {
            val conflictingSchedule = scheduleRepository
                .getScheduleForRelay(deviceId, relayIndex)
            if (conflictingSchedule?.isActive == true) {
                return@withContext Result.failure(
                    RelayLockedException(
                        "Relay locked by schedule: ${conflictingSchedule.name}"
                    )
                )
            }
        }
        
        // Create command object
        val command = Command(
            id = "${deviceId}_${System.currentTimeMillis()}",
            deviceId = deviceId,
            relay = relayIndex,
            action = action,
            source = source,
            userId = getCurrentUserId(),
            timestamp = System.currentTimeMillis(),
            executed = false,
            metadata = metadata
        )
        
        // Write to Room immediately (optimistic update)
        roomDatabase.commandDao().insert(command)
        
        // Update local relay state immediately
        val relayState = RelayState(
            deviceId = deviceId,
            relayId = relayIndex,
            state = action,
            lastToggled = System.currentTimeMillis(),
            locked = source == CommandSource.SCHEDULE
        )
        roomDatabase.relayStateDao().update(relayState)
        
        // Accumulate runtime for energy tracking
        if (action == RelayAction.ON) {
            roomDatabase.runtimeDao()
                .startTracking(deviceId, relayIndex)
        }
        
        // Send to Firebase (async, with retry on failure)
        val firebaseResult = firebaseRepository
            .sendCommand(command)
            .catch { error ->
                Log.e("applyCommand", "Firebase error: ${error.message}")
                // Command remains in queue for retry
            }
            .firstOrNull()
        
        // Post-execution logging
        logCommand(
            source = source,
            deviceId = deviceId,
            relay = relayIndex,
            action = action,
            timestamp = System.currentTimeMillis(),
            success = true
        )
        
        // Notify listeners
        _relayStateUpdated.emit(relayState)
        
        return@withContext Result.success(
            CommandResponse(
                commandId = command.id,
                relay = relayIndex,
                status = "queued",
                newState = action.toString()
            )
        )
        
    } catch (e: Exception) {
        Log.e("applyCommand", "Unexpected error", e)
        return@withContext Result.failure(e)
    }
}
\end{lstlisting}

\subsection{Android: Offline Command Queue Retry}

\begin{lstlisting}[language=Kotlin, caption=Command Queue Retry Logic, label=lst:android_queue_retry]
/**
 * Background job that processes queued commands upon connectivity restoration
 */
class CommandQueueWorker(context: Context, params: WorkerParameters) : 
    CoroutineWorker(context, params) {
    
    override suspend fun doWork(): Result = withContext(Dispatchers.IO) {
        try {
            val queuedCommands = commandRepository.getQueuedCommands()
            var successCount = 0
            var failureCount = 0
            
            for (command in queuedCommands) {
                val retryCount = command.retryCount ?: 0
                
                if (retryCount >= MAX_RETRIES) {
                    commandRepository.markCommandFailed(command.id)
                    failureCount++
                    continue
                }
                
                try {
                    // Attempt to send to Firebase
                    val response = firebaseRepository
                        .sendCommand(command)
                        .firstOrNull()
                    
                    if (response?.status == "success") {
                        commandRepository.removeFromQueue(command.id)
                        successCount++
                        Log.d("CommandQueueWorker", 
                            "Queued command executed: ${command.id}")
                    } else {
                        // Increment retry count
                        commandRepository.incrementRetry(command.id)
                        failureCount++
                    }
                    
                } catch (e: Exception) {
                    Log.w("CommandQueueWorker", "Retry failed", e)
                    commandRepository.incrementRetry(command.id)
                    failureCount++
                }
            }
            
            // Notify user of sync completion
            if (successCount > 0) {
                sendNotification(
                    "NexusFlow Queue",
                    "Synced $successCount pending commands"
                )
            }
            
            return@withContext if (failureCount == 0) 
                Result.success() 
            else 
                Result.retry()
                
        } catch (e: Exception) {
            Log.e("CommandQueueWorker", "Fatal error", e)
            return@withContext Result.retry()
        }
    }
}
\end{lstlisting}

\subsection{ESP32 Firmware: Relay Manager}

\begin{lstlisting}[language=C++, caption=ESP32 Relay Manager Module, label=lst:firmware_relay_manager]
// RelayManager.cpp - Relay control and state tracking

#include "RelayManager.h"
#include <esp_log.h>

static const char* TAG = "RelayManager";

// GPIO pin configuration
const int RELAY_PINS[RELAY_COUNT] = {19, 18, 5, 17, 16, 25};
uint32_t relay_runtime_seconds[RELAY_COUNT] = {0};
bool relay_state[RELAY_COUNT] = {false};  // false = OFF, true = ON
unsigned long relay_toggle_time[RELAY_COUNT] = {0};

void RelayManager::init() {
    ESP_LOGI(TAG, "Initializing RelayManager with %d relays", RELAY_COUNT);
    
    for (int i = 0; i < RELAY_COUNT; i++) {
        pinMode(RELAY_PINS[i], OUTPUT);
        digitalWrite(RELAY_PINS[i], HIGH);  // Active-LOW: start inactive
        relay_state[i] = false;
        relay_runtime_seconds[i] = 0;
        relay_toggle_time[i] = millis();
    }
    
    ESP_LOGI(TAG, "RelayManager initialized successfully");
}

/**
 * Apply command through unified funnel
 * All relay control paths converge here
 */
bool RelayManager::applyCommand(
    uint8_t relay_pin,
    RelayAction action,
    CommandSource source,
    uint32_t timestamp
) {
    if (relay_pin >= RELAY_COUNT) {
        ESP_LOGE(TAG, "Invalid relay pin: %d", relay_pin);
        return false;
    }
    
    // Pre-execution safety checks
    if (isRelayLocked(relay_pin)) {
        ESP_LOGW(TAG, "Relay %d locked; command rejected", relay_pin);
        return false;
    }
    
    // Execute relay control
    bool success = toggleRelay(relay_pin, action);
    if (!success) {
        ESP_LOGE(TAG, "Failed to toggle relay %d", relay_pin);
        return false;
    }
    
    // Update tracking
    updateRuntimeTracking(relay_pin);
    
    // Log command for debugging
    ESP_LOGI(TAG, "Relay %d toggled %s (source: %d, timestamp: %u)",
             relay_pin, 
             relay_state[relay_pin] ? "ON" : "OFF",
             source,
             timestamp);
    
    // Post-execution sync
    syncToFirebase(relay_pin);
    logCommand(relay_pin, action, source, timestamp);
    
    return true;
}

/**
 * Toggle relay hardware
 */
bool RelayManager::toggleRelay(uint8_t relay_pin, RelayAction action) {
    if (relay_pin >= RELAY_COUNT) return false;
    
    switch (action) {
        case RelayAction::ON:
            digitalWrite(RELAY_PINS[relay_pin], LOW);  // Active-LOW
            relay_state[relay_pin] = true;
            break;
            
        case RelayAction::OFF:
            digitalWrite(RELAY_PINS[relay_pin], HIGH);  // Inactive
            relay_state[relay_pin] = false;
            break;
            
        case RelayAction::TOGGLE:
            digitalWrite(RELAY_PINS[relay_pin], 
                        relay_state[relay_pin] ? HIGH : LOW);
            relay_state[relay_pin] = !relay_state[relay_pin];
            break;
            
        default:
            ESP_LOGE(TAG, "Unknown relay action: %d", (int)action);
            return false;
    }
    
    relay_toggle_time[relay_pin] = millis();
    return true;
}

/**
 * Accumulate relay runtime for energy tracking
 */
void RelayManager::updateRuntimeTracking(uint8_t relay_pin) {
    if (relay_pin >= RELAY_COUNT) return;
    
    static unsigned long last_update = 0;
    unsigned long now = millis();
    unsigned long delta = (now - last_update) / 1000;  // Convert to seconds
    
    if (delta > 0) {
        for (int i = 0; i < RELAY_COUNT; i++) {
            if (relay_state[i]) {  // Only accumulate when ON
                relay_runtime_seconds[i] += delta;
            }
        }
        last_update = now;
    }
}

/**
 * Get relay state
 */
bool RelayManager::getRelayState(uint8_t relay_pin) {
    if (relay_pin >= RELAY_COUNT) return false;
    return relay_state[relay_pin];
}

/**
 * Get accumulated runtime in hours
 */
float RelayManager::getRuntimeHours(uint8_t relay_pin) {
    if (relay_pin >= RELAY_COUNT) return 0.0f;
    return (float)relay_runtime_seconds[relay_pin] / 3600.0f;
}

/**
 * Check if relay is locked by schedule or automation
 */
bool RelayManager::isRelayLocked(uint8_t relay_pin) {
    // TODO: Implement conflict checking with schedule manager
    return false;
}
\end{lstlisting}

\subsection{ESP32 Firmware: Automation Engine}

\begin{lstlisting}[language=C++, caption=Automation Rule Evaluation Engine, label=lst:firmware_automation_engine]
// AutomationEngine.cpp - Rule-based device automation

#include "AutomationEngine.h"
#include "SensorManager.h"
#include "RelayManager.h"

static const char* TAG = "AutomationEngine";

#define AUTOMATION_CYCLE_INTERVAL 30000  // 30 seconds

unsigned long last_evaluation = 0;

void AutomationEngine::init() {
    ESP_LOGI(TAG, "Initializing AutomationEngine");
    last_evaluation = millis();
}

/**
 * Main automation loop - evaluates all rules every 30 seconds
 */
void AutomationEngine::update() {
    unsigned long now = millis();
    
    if ((now - last_evaluation) < AUTOMATION_CYCLE_INTERVAL) {
        return;  // Not time to evaluate yet
    }
    
    last_evaluation = now;
    
    // Get current sensor readings
    float temperature = SensorManager::getTemperature();
    float humidity = SensorManager::getHumidity();
    
    ESP_LOGI(TAG, "Evaluation cycle - Temp: %.1f deg C, Humidity: %.1f%%",
             temperature, humidity);
    
    // Evaluate all registered rules
    evaluateTemperatureThresholdRules(temperature);
    evaluateHumidityThresholdRules(humidity);
    evaluateDualSensorRules(temperature, humidity);
}

/**
 * Evaluate temperature threshold rules with hysteresis
 */
void AutomationEngine::evaluateTemperatureThresholdRules(float temp) {
    // Example: AC ON if Temp > 28 deg C with 1 deg C hysteresis
    static bool ac_rule_triggered = false;
    
    const float TEMP_THRESHOLD = 28.0f;
    const float HYSTERESIS = 1.0f;
    
    // Rising edge: trigger when temp >= threshold + hysteresis
    if (!ac_rule_triggered && temp >= (TEMP_THRESHOLD + HYSTERESIS)) {
        ESP_LOGI(TAG, "AC rule triggered: Temp %.1f >= %.1f",
                 temp, TEMP_THRESHOLD + HYSTERESIS);
        RelayManager::applyCommand(0, RelayAction::ON, 
                                   CommandSource::AUTOMATION, 
                                   millis());
        ac_rule_triggered = true;
    }
    
    // Falling edge: reset when temp <= threshold - hysteresis
    if (ac_rule_triggered && temp <= (TEMP_THRESHOLD - HYSTERESIS)) {
        ESP_LOGI(TAG, "AC rule reset: Temp %.1f <= %.1f",
                 temp, TEMP_THRESHOLD - HYSTERESIS);
        RelayManager::applyCommand(0, RelayAction::OFF, 
                                   CommandSource::AUTOMATION, 
                                   millis());
        ac_rule_triggered = false;
    }
}

/**
 * Evaluate dual-sensor rules (temperature AND humidity)
 */
void AutomationEngine::evaluateDualSensorRules(float temp, float humidity) {
    // Example: Dehumidifier ON if (Temp > 25 AND Humidity > 70)
    static bool dehumidifier_triggered = false;
    
    const float TEMP_THRESHOLD = 25.0f;
    const float HUMIDITY_THRESHOLD = 70.0f;
    const float TEMP_HYSTERESIS = 1.0f;
    const float HUMIDITY_HYSTERESIS = 3.0f;
    
    bool temp_condition = temp > (TEMP_THRESHOLD + TEMP_HYSTERESIS);
    bool humidity_condition = humidity > (HUMIDITY_THRESHOLD + HUMIDITY_HYSTERESIS);
    bool should_trigger = temp_condition && humidity_condition;
    
    if (!dehumidifier_triggered && should_trigger) {
        ESP_LOGI(TAG, "Dehumidifier ON: Temp=%.1f, Humidity=%.1f",
                 temp, humidity);
        RelayManager::applyCommand(2, RelayAction::ON, 
                                   CommandSource::AUTOMATION, 
                                   millis());
        dehumidifier_triggered = true;
    }
    
    if (dehumidifier_triggered && !should_trigger) {
        // Reset only when both conditions fall below thresholds
        bool temp_reset = temp < (TEMP_THRESHOLD - TEMP_HYSTERESIS);
        bool humidity_reset = humidity < (HUMIDITY_THRESHOLD - HUMIDITY_HYSTERESIS);
        
        if (temp_reset || humidity_reset) {
            ESP_LOGI(TAG, "Dehumidifier OFF");
            RelayManager::applyCommand(2, RelayAction::OFF, 
                                       CommandSource::AUTOMATION, 
                                       millis());
            dehumidifier_triggered = false;
        }
    }
}

/**
 * Log rule evaluation for debugging
 */
void AutomationEngine::logRuleExecution(
    const char* rule_name,
    bool executed,
    int relay_pin,
    RelayAction action
) {
    if (executed) {
        ESP_LOGI(TAG, "Rule '%s' executed: Relay %d -> %s",
                 rule_name, relay_pin,
                 action == RelayAction::ON ? "ON" : "OFF");
    }
}
\end{lstlisting}

\subsection{Android: Presence Heartbeat}

\begin{lstlisting}[language=Kotlin, caption=Presence Sync Mechanism, label=lst:android_presence]
/**
 * Manages user presence for adaptive Firebase sync intervals
 */
class PresenceManager(
    private val firebaseDatabase: FirebaseDatabase,
    private val userId: String
) {
    
    private val presenceReference = 
        firebaseDatabase.reference.child("presence/$userId")
    
    /**
     * Write presence heartbeat when app becomes active
     */
    fun onAppForeground(deviceId: String) {
        val presenceData = mapOf(
            "lastActive" to System.currentTimeMillis(),
            "appVersion" to BuildConfig.VERSION_NAME,
            "userId" to userId
        )
        
        presenceReference.child(deviceId).setValue(presenceData)
            .addOnSuccessListener {
                Log.d("PresenceManager", "Presence updated: $deviceId")
            }
            .addOnFailureListener { error ->
                Log.w("PresenceManager", "Presence update failed", error)
            }
        
        // Set up onDisconnect cleanup
        presenceReference.child(deviceId)
            .onDisconnect()
            .removeValue()
            .addOnFailureListener { error ->
                Log.w("PresenceManager", "onDisconnect setup failed", error)
            }
    }
    
    /**
     * Called from lifecycle when app enters background
     * Firebase onDisconnect listener will clean up automatically
     */
    fun onAppBackground() {
        // Explicit cleanup not needed; onDisconnect handles it
        Log.d("PresenceManager", "App backgrounded; onDisconnect active")
    }
}

/**
 * Adjust Firebase sync interval based on presence state
 */
class PresenceAwareSyncManager(
    private val firebaseDatabase: FirebaseDatabase,
    private val deviceId: String
) {
    
    private val presenceReference = firebaseDatabase.reference
        .child("presence/$deviceId")
    
    fun startMonitoringPresence() {
        presenceReference.addValueEventListener(
            object : ValueEventListener {
                override fun onDataChange(snapshot: DataSnapshot) {
                    val presence = snapshot.value as? Map<*, *>
                    
                    if (presence != null && presence.isNotEmpty()) {
                        // User is active - use 1-minute sync interval
                        setSyncInterval(SYNC_INTERVAL_ACTIVE_MS)
                        Log.d("PresenceSyncManager", "User active: 1-min sync")
                    } else {
                        // User is idle - use 15-minute sync interval
                        setSyncInterval(SYNC_INTERVAL_IDLE_MS)
                        Log.d("PresenceSyncManager", "User idle: 15-min sync")
                    }
                }
                
                override fun onCancelled(error: DatabaseError) {
                    Log.e("PresenceSyncManager", 
                        "Presence monitoring error", error.toException())
                }
            }
        )
    }
    
    private fun setSyncInterval(intervalMs: Long) {
        // Update Firebase listener with new interval
        // Implementation depends on sync mechanism
    }
    
    companion object {
        private const val SYNC_INTERVAL_ACTIVE_MS = 60_000   // 1 minute
        private const val SYNC_INTERVAL_IDLE_MS = 900_000    // 15 minutes
    }
}
\end{lstlisting}

\section{Compilation and Deployment}

\subsection{Android Build and Installation}

\begin{lstlisting}[language=bash, caption=Android Build Commands, label=lst:android_build]
# Build debug APK
./gradlew assembleDebug

# Install on connected device
./gradlew installDebug

# Build release APK (requires signing configuration)
./gradlew assembleRelease

# Run unit tests
./gradlew testDebugUnitTest

# Run instrumented tests on device
./gradlew connectedAndroidTest
\end{lstlisting}

\subsection{ESP32 Firmware Build}

\begin{lstlisting}[language=bash, caption=ESP32 Build and Flash, label=lst:esp32_build]
# Build firmware
idf.py build

# Flash to device
idf.py -p /dev/ttyUSB0 flash

# Monitor serial output
idf.py -p /dev/ttyUSB0 monitor

# Build and flash in one command
idf.py -p /dev/ttyUSB0 flash monitor

# Full clean rebuild
idf.py fullclean build
\end{lstlisting}

\section{Testing and Debugging}

\subsection{Android Debugging}

\begin{lstlisting}[language=bash, caption=Android Debug Tools, label=lst:android_debug]
# List connected devices
adb devices

# View live logs
adb logcat | grep "NexusFlow"

# Dump Room database
adb shell "sqlite3 /data/data/com.nexusflow.app/databases/nexusflow_db.db"

# Monitor network traffic (requires tcpdump)
adb shell "su 0 tcpdump -i any -w /sdcard/capture.pcap"

# Inspect Firebase RTDB locally
# Use Firebase Emulator Suite
firebase emulators:start
\end{lstlisting}

\subsection{ESP32 Serial Debugging}

\begin{lstlisting}[language=bash, caption=ESP32 Serial Monitoring, label=lst:esp32_debug]
# Monitor with timestamps
idf.py monitor --timestamps

# Log specific module
idf.py monitor --tag RelayManager

# Export logs to file
idf.py monitor > esp32_logs.txt 2>&1
\end{lstlisting}

These code snippets represent the core functional components of NexusFlow, demonstrating best practices in mobile development, embedded systems, and IoT integration.
